\section{Model description: Model X}

\textit{Note: This section defines the general two-class system equipped with
\emph{both} jockeying ($\gamma_i$) and abandonments ($\theta_i$). We refer to this
full model as \emph{Model X}. The simpler variants studied elsewhere in this
document are recovered as special cases by zeroing the appropriate parameters, as
summarised below.}

\begin{table}[H]
    \centering
    \begin{tabular}{cccl}
        \hline
        Model & Jockeying & Abandonment & Parameter restriction \\
        \hline
        X  & both ($1\!\leftrightarrow\!2$) & both ($1,2$)       & $\gamma_i \geq 0,\ \theta_i \geq 0$ \quad (general) \\
        A  & none                            & none               & $\gamma_1=\gamma_2=0,\ \theta_1=\theta_2=0$ \\
        B  & both ($1\!\leftrightarrow\!2$) & none               & $\gamma_i \geq 0,\ \theta_1=\theta_2=0$ \\
        B2 & $1\!\to\!2$ only                & none               & $\gamma_1 \geq 0,\ \gamma_2=0,\ \theta_1=\theta_2=0$ \\
        C2 & none                            & class $1$ only     & $\theta_1 \geq 0,\ \theta_2=0,\ \gamma_1=\gamma_2=0$ \\
        \hline
    \end{tabular}
    \caption{Nested family of models. Model X is the parent; Models A, B, B2 and C2
    are obtained by restricting the jockeying directions and the abandoning classes.}
    \label{tab:model_family}
\end{table}

The model discussed in this section is an $M/M/1$ queueing system with $2$ priority
classes, where customers are served on a First-Come, First-Served (FCFS) basis
within their own class and those coming from the priority class (Class-$1$) have
non-preemptive priority over those of the non-priority class (Class-$2$); that is,
when a Class-$2$ customer is in service, that service will not be interrupted by a
potential arrival of a Class-$1$ customer. Beyond the basic dynamics, Model X
incorporates two distinct queue-leaving mechanisms. First, \emph{jockeying}: a
waiting customer may abandon its own queue and \emph{join the other queue}, at a
per-customer rate $\gamma_1$ for the direction $1\to2$ and $\gamma_2$ for the
direction $2\to1$, so the total number of customers is preserved. Second,
\emph{abandonment} (reneging): a waiting customer may leave the \emph{system
entirely} at a per-customer rate $\theta_i$, so the total number of customers
decreases by one. Concretely, each Class-$i$ customer that is waiting in line
(i.e.\ not the one in service) holds an independent $\mathrm{Exp}(\theta_i)$
patience clock and abandons the system when it expires; consequently the aggregate
Class-$i$ abandonment rate in a state with $n_i$ waiting customers is the linear
term $\theta_i n_i$ that appears in the balance equations. This is precisely the
linear-reneging mechanism of the $M/M/1+M$ queue introduced in the Preliminaries.

\begin{figure}[htbp]
    \centering
    \input{figures/queue_diagram.tikz}
    \caption{Queueing system with two parallel queues.}
    \label{fig:queue_model}
\end{figure}

This model, being an $M/M/1$ queueing model according to Kendall's notation,
features arrivals driven by a Poisson process with arrival rates $\lambda_i$ for
Class-$i$, exponentially distributed service times with rate $\mu$, and a service
capacity of $1$ customer at a time. In addition, a waiting customer of Class-$1$
jockeys to the class-$2$ queue at per-customer rate $\gamma_1$ and a waiting
customer of Class-$2$ jockeys to the class-$1$ queue at per-customer rate
$\gamma_2$ (transfers that conserve the total), while a waiting customer of
Class-$i$ abandons the system at per-customer rate $\theta_i$ (a true departure
that reduces the total).

In order to define the stochastic processes that drive this queueing system, it is
necessary to introduce the following random variables:
\begin{itemize}
    \item $\mathcal{B}:$ binary indicator for the server being idle or busy, $\mathcal{B}\in \{0,1\}$;
    \item $N_1, N_2:$ number of customers of class $1$ and $2$, respectively, in their respective queues;
    \item $N:$ total number of customers in the queues.
\end{itemize}
Note that the variables $N_1$, $N_2$ and $N$ count the number of customers in the
queues, as opposed to the customers in the system. This definition is convenient
given that we have a common service rate for both priority classes, and hence the
states of the system are equivalent regardless of the priority class the customer
in service comes from. Since having customers in line directly implies that someone
is in service, we will denote the idle state as $(0)$ and the rest of the states by
the number of people in each queue as the tuple $(n_1,n_2)$, where $n_i$ denotes the
number of customers in the queue for class-$i$. This is the most conventional way to
define the state space, namely $S$, in such a queueing system. However, we can also
look at the total number of customers in the queues. We can define an alternative
state space $\widetilde{S}$ using the idle state $(0)$ and the pair $(n_2,n)$ —we
could also choose $(n_1,n)$—, which gives that we have $(n-n_2)$ Class-$1$ customers
in the priority queue. Note that, since $n$ is the total number of customers, we
need the constraint $n \geq n_2$.
Therefore, we define our state spaces as follows:
\begin{align}
    \begin{split}
        &S := \{(0)\} \cup \{ (n_1,n_2) \mid n_1,n_2 \geq 0 \},
        \label{eq:state_space_s}
    \end{split} \\
    \begin{split}
        &\widetilde{S} := \{(0)\} \cup \{ (n_2,n) \mid n \geq n_2 \geq 0 \}.
        \label{eq:state_space_stilde}
    \end{split}
\end{align}

For these state spaces, we define the limiting —as opposed to time-dependent—
probabilities of finding the system in that state in the long run as follows:
\begin{itemize}
    \item $\pi_0$ for the idle state $(0)$, which belongs to both spaces $S$ and $\widetilde{S}$;
    \item $\pi(n_1,n_2)$ for states $(n_1,n_2) \in S$, where $n_1,n_2 \geq 0$;
    \item $\widetilde{\pi}(n_2,n)$ for states $(n_2,n) \in \widetilde{S}$, where $n \geq n_2 \geq 0$.
\end{itemize}

Regarding the parameters of the system, the arrival process for each class follows a
Poisson process (i.e., the interarrival times are exponentially distributed). By the
properties of Poisson processes, the total number of arrivals will also form a
Poisson process. Furthermore, we have a common service rate for exponentially
distributed service times. We can organize our notation into class-specific
parameters (alongside our common service rate) and class-blind totals.
For the class-specific metrics and service rate, we define:
\begin{itemize}
    \item $\lambda_1,\lambda_2$: arrival rates of classes $1$ and $2$, respectively;
    \item $\mu$: service rate of a customer, regardless of their priority class;
    \item $\rho_i=\frac{\lambda_i}{\mu}$: system load of class $i$, for $i\in\{1,2\}$;
    \item $\gamma_1$: per-customer rate at which a waiting Class-$1$ customer jockeys
          from the class-$1$ queue to the class-$2$ queue ($1\to2$);
    \item $\gamma_2$: per-customer rate at which a waiting Class-$2$ customer jockeys
          from the class-$2$ queue to the class-$1$ queue ($2\to1$); a jockeying
          event in either direction transfers a waiting customer between queues and
          \emph{conserves} the total $n=n_1+n_2$ (Model B2 sets $\gamma_2=0$);
    \item $\theta_1,\theta_2$: per-customer abandonment (reneging) rates from the
          class-$1$ and class-$2$ queues, respectively; an abandonment is a true
          departure from the system and \emph{reduces} the total number of
          customers by one. Each waiting Class-$i$ customer reneges independently
          after an $\mathrm{Exp}(\theta_i)$ patience time, so the aggregate
          Class-$i$ abandonment rate in state $(n_1,n_2)$ is $\theta_i n_i$
          (Model C2 sets $\theta_2=0$).
\end{itemize}
For the aggregate system parameters, we define:
\begin{itemize}
    \item $\lambda=\lambda_1+\lambda_2$: total arrival rate;
    \item $\rho=\frac{\lambda}{\mu}$: total system load.
\end{itemize}

\textbf{Stability.}
In the absence of abandonment ($\theta_1=\theta_2=0$, i.e.\ Models A, B and B2), the
server must keep up on average with the offered workload, and it suffices to require
\begin{equation*}
    \rho=\frac{\lambda_1+\lambda_2}{\mu}<1 .
\end{equation*}
Jockeying does not affect this condition, since it merely relocates customers
without changing their total number. Abandonment changes the picture: when
\emph{both} $\theta_1>0$ and $\theta_2>0$ (the general Model X), the state-dependent
departure rate $\mu+\theta_1 n_1+\theta_2 n_2$ grows without bound in every
direction, so the chain is positive recurrent for \emph{all} $\lambda,\mu>0$. When
only one class reneges —as in Model C2, where $\theta_1>0$ but $\theta_2=0$—
abandonment bounds only the priority queue, and stability still imposes a condition
on the non-abandoning class; this is treated in the dedicated analysis of that
model.

\textbf{Computation of $\pi_0$ and $\pi(0,0)$.}
A relation that holds for Model X \emph{unconditionally} —for any
$\gamma_i,\theta_i\geq 0$— follows from the cut between the idle state and the set of
busy states. The only transition into $(0)$ is a service completion from the
single-customer state $(0,0)$ at rate $\mu$, and the only transition out of $(0)$ is
an arrival at rate $\lambda=\lambda_1+\lambda_2$; this is exactly the idle balance
equation $\lambda\pi_0=\mu\,\pi(0,0)$. Consequently,
\begin{equation}
    \pi(0,0) \;=\; \frac{\lambda}{\mu}\,\pi_0 \;=\; \rho\,\pi_0 .
    \label{eq:pi00_general}
\end{equation}
Note that neither jockeying nor abandonment can occur from $(0,0)$ (there are no
waiting customers), so~\eqref{eq:pi00_general} is unaffected by $\gamma_i$ and
$\theta_i$. What \emph{does} depend on the abandonment parameters is the value of
$\pi_0$ itself, as we now discuss.

\emph{Models without abandonment (A, B, B2).} When $\theta_1=\theta_2=0$, jockeying
conserves the total number of customers, so the total-in-system process is a
constant death-rate birth--death chain with up-rate $\lambda$ and down-rate $\mu$
—that is, an ordinary $M/M/1$ queue with parameters $\lambda$ and $\mu$. Its idle
probability is the familiar
\begin{equation}
    \pi_0 = 1-\rho = 1-\frac{\lambda_1+\lambda_2}{\mu},
    \qquad
    \pi(0,0)=\rho\,\pi_0=(1-\rho)\rho .
    \label{eq:pi0_noaband}
\end{equation}

\emph{Models with abandonment (X, C2).} As soon as a true abandonment is present, the
total number of customers is no longer conserved: the departure rate out of a busy
state is $\mu+\theta_1 n_1+\theta_2 n_2$, which depends on the \emph{split}
$(n_1,n_2)$ and not merely on the total $n=n_1+n_2$. Indeed, the states $(n,0)$ and
$(0,n)$ carry departure rates $\mu+\theta_1 n$ and $\mu+\theta_2 n$, which differ
whenever $\theta_1\neq\theta_2$ —in particular for Model C2, where $\theta_2=0$.
Consequently the total-in-system process is not lumpable to a one-dimensional
birth--death chain, and $\pi_0$ admits no closed form obtained from the marginal
alone. Instead, $\pi_0$ is determined jointly with the full two-dimensional
stationary distribution through the normalisation $P(1,1)=1-\pi_0$ of the bivariate
generating function, as carried out in the analysis sections below —and, for the
single-class-abandonment case $\theta_2=0$, in the dedicated analysis of Model C2.

In all cases, the mass on the busy states is
\begin{equation}
    \sum_{n_1=0}^{\infty}\sum_{n_2=0}^{\infty}\pi(n_1,n_2) = 1-\pi_0,
    \label{eq:busy_mass}
\end{equation}
which equals $\rho$ only in the no-abandonment case~\eqref{eq:pi0_noaband}; for
Models X and C2 one has $1-\pi_0\neq\rho$ in general. This conditioning on a busy
server will be relevant in our derivations, as we will look at the queue length
conditioned on the server being busy.

\subsection{Balance equations}

\subsubsection{Balance equations for state space $S$} \label{sec:n1n2_balance_equations}

Diagram for the combined model. Dashed arrows correspond to jockeying and dotted arrows to abandonments; both belong only to the generalized model, while the solid arrows are common to all variants:
\begin{figure}[H]
  \centering
  \resizebox{0.8\textwidth}{!}{%
    \begin{tikzpicture}[>=Stealth, line cap=round, line join=round]
  % --- global axes ---
  \draw[->] (0,0) -- (8,0) node[below right] {$n_1$};
  \draw[->] (0,0) -- (0,8) node[above left] {$n_2$};
  
  % Origin label
  \node[above right, xshift=2pt, yshift=2pt] at (0,0) {$(0,0)$};

  % --- Entry transitions from the origin ---
  \draw[->] (0,0) -- ++(0:1.2)  node[below] {$\lambda_1$};
  \draw[->] (0,0) -- ++(90:1.2) node[right] {$\lambda_2$};

  % --- boundary state (n1,0) ---
  \coordinate (A) at (6,0);
  \node[below] at (A) {$(n_1,0)$};
  \draw[->] (A) -- ++(0:1.2)   node[below] {$\lambda_1$}; 
  \draw[->] (A) -- ++(90:1.2)  node[right] {$\lambda_2$}; 
  \draw[->] (A) -- ++(180:1.2) node[below] {$\mu$};      
  \draw[->, dashed] (A) -- ++(135:1.2) node[left]  {$\gamma_1 n_1$};
  % Abandonment of class-1: parallel dotted arrow above mu
  \draw[->, dotted] (6, 0.3) -- (4.8, 0.3) node[above] {$\theta_1 n_1$};

  % --- interior state (n1,n2) ---
  \coordinate (B) at (6,4);
  \node[below left] at (B) {$(n_1,n_2)$};
  \draw[->] (B) -- ++(0:1.2)   node[below] {$\lambda_1$};
  \draw[->] (B) -- ++(90:1.2)  node[right] {$\lambda_2$};
  \draw[->] (B) -- ++(180:1.2) node[left]  {$\mu$};
  \draw[->, dashed] (B) -- ++(135:1.2) node[left]  {$\gamma_1 n_1$};
  \draw[->, dashed] (B) -- ++(315:1.2) node[right] {$\gamma_2 n_2$};
  % Abandonments: parallel dotted arrows
  \draw[->, dotted] (6, 4.3) -- (4.8, 4.3) node[above] {$\theta_1 n_1$};
  \draw[->, dotted] (B)      -- ++(270:1.2) node[left]  {$\theta_2 n_2$};

  % --- boundary state (0,n2) ---
  \coordinate (C) at (0,4);
  \node[left] at (C) {$(0,n_2)$};
  \draw[->] (C) -- ++(0:1.2)   node[below] {$\lambda_1$};
  \draw[->] (C) -- ++(90:1.2)  node[right] {$\lambda_2$};
  \draw[->] (C) -- ++(-90:1.2) node[left]  {$\mu$};  
  \draw[->, dashed] (C) -- ++(315:1.2) node[right] {$\gamma_2 n_2$};
  % Abandonment of class-2: parallel dotted arrow to the right of mu
  \draw[->, dotted] (0.3, 4) -- (0.3, 2.8) node[right] {$\theta_2 n_2$};

  % --- "idle" arcs ---
  \coordinate (idle) at (-1.25,-1.15);
  \node[below] at (idle) {$\text{idle}$};
  \draw[->] (idle) to[bend left=45] node[midway, above left] {$\lambda_1+\lambda_2$} (0,0);
  \draw[->] (0,0) to[bend left=45] node[midway, below right] {$\mu$} (idle);
\end{tikzpicture}
  }
  \caption{State transition diagram on state space $S$ for the model with jockeying ($\gamma_i n_i$, dashed) and abandonments ($\theta_i n_i$, dotted).}
  \label{fig:diagram-gen}
\end{figure}

\begin{align}
    \begin{split}
        &\left[\lambda_1 + \lambda_2 + \mu + (\gamma_1 + \theta_1) n_1 + (\gamma_2 + \theta_2) n_2\right] \pi(n_1,n_2) \\
        & \quad = \quad \mu \pi(n_1+1,n_2) \\
        & \quad \quad + \lambda_1 \pi(n_1-1,n_2) \\
        & \quad \quad + \lambda_2 \pi(n_1,n_2-1) \\
        & \quad \quad + \gamma_1 (n_1+1) \pi(n_1+1,n_2-1) \\
        & \quad \quad + \gamma_2 (n_2+1) \pi(n_1-1,n_2+1) \\
        & \quad \quad + \theta_1 (n_1+1) \pi(n_1+1,n_2) \\
        & \quad \quad + \theta_2 (n_2+1) \pi(n_1,n_2+1)
    \end{split}
    && \text{for} \quad n_1 \geq 1, n_2 \geq 1 \label{eq:gen_balance_interior}\\
    \begin{split}
        & \left[ \lambda_1 + \lambda_2 + \mu +  (\gamma_2 + \theta_2) n_2 \right] \pi(0,n_2) \\
        & \quad = \quad \lambda_2 \pi(0,n_2-1) \\
        & \quad \quad + \mu \pi(1,n_2) \\
        & \quad \quad + \mu \pi(0,n_2+1) \\
        & \quad \quad + \gamma_1 \pi(1, n_2-1) \\
        & \quad \quad + \theta_1 \pi(1, n_2) \\
        & \quad \quad + \theta_2(n_2+1) \pi(0,n_2+1)
    \end{split}
    && \text{for} \quad n_1=0, n_2 \geq 1 \label{eq:gen_balance_y-boundary}\\
    \begin{split}
        & \left[ \lambda_1 + \lambda_2 + \mu + (\gamma_1 + \theta_1) n_1 \right] \pi(n_1,0) \\
        & \quad = \quad \lambda_1 \pi(n_1-1,0) \\
        & \quad \quad + \mu \pi(n_1+1, 0) \\
        & \quad \quad + \gamma_2 \pi(n_1-1, 1) \\
        & \quad \quad + \theta_1(n_1+1) \pi(n_1+1, 0) \\
        & \quad \quad + \theta_2 \pi(n_1,1)
    \end{split}
    && \text{for} \quad n_1 \geq 1, n_2 = 0 \label{eq:gen_balance_x-boundary}\\
    \begin{split}
        \left[\lambda_1+\lambda_2\right]\pi(0,0) = (\mu+\theta_1)\pi(1,0) + (\mu+\theta_2)\pi(0,1)
    \end{split} && \label{eq:gen_balance_zero}\\
    \begin{split}
        (\lambda_1+\lambda_2) \pi_0 = \mu \pi(0,0)
    \end{split} && \label{eq:gen_balance_idle}
\end{align}

\subsubsection{Balance equations for state space $\widetilde{S}$}

We introduce a new state space $\widetilde{S}$:
\begin{equation}
    \widetilde{S} := \big\{ (0) \cup \{ (n_2,n) \mid n \geq n_2 \geq 0\}\big\}.
\end{equation}

For the sake of readability, we will denote the limiting probabilities associated to this state space with $\tpi(n_2,n)$, for state $(n_2,n)\in \widetilde{S}$.

\begin{figure}[H]
  \centering
  \resizebox{0.8\textwidth}{!}{%
    \input{figures/diagram_jockeying_and_abandonments.tikz}
  }
  \caption{State transition diagram on state space $\widetilde{S}$.}
\end{figure}

The balance equations, after combining the equations where $\tpi_0$ shows up, are:
\begin{align}
    \begin{split}
        & \left[ \lambda_1 + \lambda_2 + \mu + (\gamma_1+\theta_1)(n-n_2) + (\gamma_2+\theta_2) n_2 \right] \tpi(n_2,n) \\
        & \quad = \quad \lambda_1 \tpi(n_2,n-1) \\
        & \quad \quad + \lambda_2 \tpi(n_2-1,n-1) \\
        & \quad \quad + \mu \tpi(n_2, n+1) \\
        & \quad \quad + \gamma_1 (n-n_2+1) \tpi(n_2-1, n) \\
        & \quad \quad + \gamma_2 (n_2+1) \tpi(n_2+1,n) \\
        & \quad \quad + \theta_1 (n+1-n_2) \tpi(n_2, n+1) \\
        & \quad \quad + \theta_2 (n_2+1) \tpi(n_2+1, n+1)
    \end{split}
    && \text{for} \quad 1 \leq n_2 \leq n-1, \; n \geq 2 \label{eq:tilde_balance_interior} \\
    \begin{split}
        & \left[ \lambda_1 + \lambda_2 + \mu + (\gamma_2+\theta_2) n \right] \tpi(n,n) \\
        & \quad = \quad \lambda_2 \tpi(n-1,n-1) \\
        & \quad \quad + \mu \tpi(n+1,n+1) \\
        & \quad \quad + \mu \tpi(n,n+1) \\
        & \quad \quad + \gamma_1 \tpi(n-1, n) \\
        & \quad \quad + \theta_1 \tpi(n, n+1) \\
        & \quad \quad + \theta_2 (n+1) \tpi(n+1, n+1)
    \end{split}
    && \text{for} \quad n_2=n, \; n \geq 1 \label{eq:tilde_balance_diagonal} \\
    \begin{split}
        & \left[ \lambda_1 + \lambda_2 + \mu + (\gamma_1+\theta_1) n \right] \tpi(0,n) \\
        & \quad = \quad \lambda_1 \tpi(0,n-1) \\
        & \quad \quad + \mu \tpi(0, n+1) \\
        & \quad \quad + \gamma_2 \tpi(1,n) \\
        & \quad \quad + \theta_1 (n+1) \tpi(0, n+1) \\
        & \quad \quad + \theta_2 \tpi(1, n+1)
    \end{split}
    && \text{for} \quad n_2=0, \; n \geq 1 \label{eq:tilde_balance_y-boundary} \\
    \begin{split}
        & \left[ \lambda_1 + \lambda_2 \right] \tpi(0,0) \\
        & \quad = \quad (\mu + \theta_1) \tpi(0,1) + (\mu + \theta_2) \tpi(1,1)
    \end{split}
    && \label{eq:tilde_balance_zero} \\
    \begin{split}
        (\lambda_1+\lambda_2) \tpi_0 = \mu \tpi(0,0)
    \end{split}
    && \label{eq:tilde_balance_idle}
\end{align}


\subsection{Derivation of equations for (Partial) Probability Generating functions}

In order to study the limiting behaviour of this two-class system, we define the bivariate Probability Generating Function (PGF) of the joint number of customers in the queue for each class, $N_1$ and $N_2$:
\begin{equation}
    P(x,y) = \mathbb{E}[x^{N_1} y^{N_2}] = \sum_{n_1=0}^{\infty} \sum_{n_2=0}^{\infty} \pi(n_1,n_2) x^{n_1} y^{n_2}.
\end{equation}

The derivation that will be carried out here includes the boundary terms, where at least one of the two queues is empty:
\begin{align}
    P_x(x) &= P(x,0) = \sum_{n_1=0}^{\infty} \pi(n_1,0) x^{n_1} \\
    P_y(y) &= P(0,y) = \sum_{n_2=0}^{\infty} \pi(0,n_2) y^{n_2}.
\end{align}

Note that both of these terms include the state where the server is busy but both queues are empty, which is represented by $\pi(0,0)$, while the state where the entire system is empty (the idle state) is treated separately.


Another term that arises during the application of the balance equations characterizes the boundary dynamics when the Class-$1$ contains exactly one customer:
\begin{equation}
    P_1(y) = \sum_{n_2=1}^{\infty} \pi(1,n_2) y^{n_2}.
\end{equation}
This term serves as an intermediate step; as the derivation proceeds, it will be possible to express it in terms of the primary boundary functions $P_x(x)$ and $P_y(y)$.

In the same fashion, we introduce the concept of the \textbf{Partial Probability Generating Function (PPGF)}. This technique, notably utilized by Cohen \cite{Cohen1950}, consists of transforming only a subset of the random variables while maintaining the discrete indices of the remainder. In the context of a priority queue, this method is particularly advantageous: it allows us to exploit the well-known properties of the priority class (often a standard $M/M/1$ system) and focus the analytical effort on the non-trivial PGF of the non-priority class.

While Cohen originally developed this framework for multi-server systems, we apply it here to the single-server case. To avoid the potential confusion stemming from mid-20th-century notation, we provide a modern formulation consistent with the variables defined above.

We define the PPGF with respect to the second variable $N_2$, conditioned on the first queue being in state $n_1$, as:
\begin{equation}
    \widetilde{P}(n_1, y) = \mathbb{E} \left[y^{N_2} \mathbb{1}_{\{ N_1=n_1\}} \right] = \sum_{n_2=0}^{\infty} \pi(n_1,n_2) y^{n_2}.
\end{equation}

Similarly, the transform with respect to the first variable $N_1$, while maintaining the discrete state $n_2$ for the second queue, is defined as:
\begin{equation}
    \widetilde{P}(x, n_2) = \mathbb{E} \left[x^{N_1} \mathbb{1}_{\{ N_2=n_2\}} \right] = \sum_{n_1=0}^{\infty} \pi(n_1,n_2) x^{n_1}.
\end{equation}

These partial transforms are linked to the joint PGF by the following relations:
\begin{equation}
    P(x,y) = \sum_{n_1=0}^{\infty} \widetilde{P}(n_1, y) x^{n_1} = \sum_{n_2=0}^{\infty} \widetilde{P}(x, n_2) y^{n_2}.
\end{equation}


\subsubsection{Equation for PGF in state space $S$}

\begin{lemma}[Fundamental equation for the PGF on $S$]\label{lem:gen_fundamental_equation}
Consider the two-class non-preemptive priority queue on the state space
$S=\{(n_1,n_2):n_1,n_2\ge 0\}$, with Poisson arrival rates $\lambda_1,\lambda_2$,
exponential service rate $\mu$, jockeying rates $\gamma_1,\gamma_2$ and abandonment
rates $\theta_1,\theta_2$, where $n_1$ and $n_2$ count the class-$1$ and class-$2$
customers \emph{waiting} in the queues. Assume the associated CTMC is positive
recurrent, so that the stationary distribution $\{\pi(n_1,n_2)\}$ exists and the
joint PGF
\[
    P(x,y)=\sum_{n_1=0}^{\infty}\sum_{n_2=0}^{\infty}\pi(n_1,n_2)\,x^{n_1}y^{n_2}
\]
is analytic on the closed unit polydisc $|x|\le 1,\ |y|\le 1$. Then $P(x,y)$,
together with the boundary function $P_y(y)=P(0,y)=\sum_{n_2=0}^{\infty}\pi(0,n_2)y^{n_2}$,
satisfies the first-order linear partial differential equation
\begin{equation}
\begin{split}
    & \left[ (\lambda_1+\lambda_2+\mu)xy - \mu y - \lambda_1 x^2 y - \lambda_2 x y^2 \right] P(x,y) \\
    & + xy \left[ \gamma_1(x-y) + \theta_1(x-1) \right] \frac{\partial}{\partial x} P(x,y) \\
    & + xy \left[ \gamma_2(y-x) + \theta_2(y-1) \right] \frac{\partial}{\partial y} P(x,y) \\
    & \quad = \quad \mu(x-y) P_y(y) - \mu x(1-y)\,\pi(0,0).
\end{split}
\label{eq:gen_fundamental_equation}
\end{equation}
\end{lemma}

\begin{proof}

    The balance equations are as defined in Equations~\eqref{eq:gen_balance_interior}, \eqref{eq:gen_balance_y-boundary}, \eqref{eq:gen_balance_x-boundary}, \eqref{eq:gen_balance_zero} and \eqref{eq:gen_balance_idle}. We follow the standard procedure for finding PGF's: multiply every balance equation by $x^{n_1}y^{n_2}$ and add them up across their domain.
    
    Hence, the interior balance equation (Equation~\eqref{eq:gen_balance_interior}) yields the following expression:
    \begin{align}
        \begin{split}
            & (\lambda_1 + \lambda_2 + \mu) \sum_{n_1=1}^{\infty} \sum_{n_2=1}^{\infty} \pi(n_1,n_2) x^{n_1}y^{n_2} \\
            &+ \gamma_1 \sum_{n_1=1}^{\infty} \sum_{n_2=1}^{\infty} n_1 \pi(n_1,n_2) x^{n_1}y^{n_2} \\
            &+ \gamma_2 \sum_{n_1=1}^{\infty} \sum_{n_2=1}^{\infty} n_2 \pi(n_1,n_2) x^{n_1}y^{n_2} \\
            &+ \theta_1 \sum_{n_1=1}^{\infty} \sum_{n_2=1}^{\infty} n_1 \pi(n_1,n_2) x^{n_1}y^{n_2} \\
            &+ \theta_2 \sum_{n_1=1}^{\infty} \sum_{n_2=1}^{\infty} n_2 \pi(n_1,n_2) x^{n_1}y^{n_2} \\
            &\quad = \quad \mu \sum_{n_1=1}^{\infty} \sum_{n_2=1}^{\infty} \pi(n_1+1,n_2)x^{n_1}y^{n_2} \\
            &\quad\quad  + \lambda_1 \sum_{n_1=1}^{\infty} \sum_{n_2=1}^{\infty} \pi(n_1-1,n_2)x^{n_1}y^{n_2} \\
            &\quad\quad + \lambda_2 \sum_{n_1=1}^{\infty} \sum_{n_2=1}^{\infty} \pi(n_1,n_2-1)x^{n_1}y^{n_2} \\
            &\quad\quad + \gamma_1 \sum_{n_1=1}^{\infty} \sum_{n_2=1}^{\infty} (n_1+1)\pi(n_1+1,n_2-1)x^{n_1}y^{n_2} \\
            &\quad\quad + \gamma_2 \sum_{n_1=1}^{\infty} \sum_{n_2=1}^{\infty} (n_2+1)\pi(n_1-1,n_2+1)x^{n_1}y^{n_2} \\
            &\quad\quad + \theta_1 \sum_{n_1=1}^{\infty} \sum_{n_2=1}^{\infty} (n_1+1)\pi(n_1+1,n_2)x^{n_1}y^{n_2} \\
            &\quad\quad + \theta_2 \sum_{n_1=1}^{\infty} \sum_{n_2=1}^{\infty} (n_2+1)\pi(n_1,n_2+1)x^{n_1}y^{n_2}
        \end{split} \\
        \begin{split}
            & LHS_1 + LHS_2 + LHS_3 + LHS_4 + LHS_5 \\
            & \quad = \quad RHS_1 + RHS_2 + RHS_3 + RHS_4 + RHS_5 + RHS_6 + RHS_7
        \end{split}
    \end{align}
    
    We solve the parts individually:
    \begin{align}
        \begin{split}
            & LHS_1 \\
            & \quad = \quad (\lambda_1 + \lambda_2 + \mu) \sum_{n_1=1}^{\infty} \sum_{n_2=1}^{\infty} \pi(n_1,n_2)x^{n_1}y^{n_2} \\
            & \quad = \quad (\lambda_1 + \lambda_2 + \mu) \sum_{n_1=1}^{\infty} \left[ \sum_{n_2=0}^{\infty} \pi(n_1,n_2) x^{n_1}y^{n_2} - \pi(n_1,0) x^{n_1} y^0 \right] \\
            & \quad = \quad (\lambda_1 + \lambda_2 + \mu) \left[
                \sum_{n_1=1}^{\infty} \sum_{n_2=0}^{\infty} \pi(n_1,n_2) x^{n_1}y^{n_2}
                - \sum_{n_1=1}^{\infty} \pi(n_1,0) x^{n_1}
            \right] \\
            & \quad = \quad (\lambda_1 + \lambda_2 + \mu) \left[
                \sum_{n_1=0}^{\infty} \sum_{n_2=0}^{\infty} \pi(n_1,n_2) x^{n_1}y^{n_2}
                - \sum_{n_2=0}^{\infty} \pi(0,n_2) y^{n_2}
            \right] \\
            & \quad \quad - \; (\lambda_1 + \lambda_2 + \mu) \left[
                \sum_{n_1=0}^{\infty} \pi(n_1,0) x^{n_1} - \pi(0,0)
            \right] \\
            & \quad = \quad (\lambda_1 + \lambda_2 + \mu) \left[ P(x,y) - P_y(y) - P_x(x) + \pi(0,0) \right]
        \end{split} \\
        \begin{split}
            & LHS_2 \\
            & \quad = \quad \gamma_1 \sum_{n_1=1}^{\infty} \sum_{n_2=1}^{\infty} n_1 \pi(n_1,n_2) x^{n_1} y^{n_2} \\
            & \quad = \quad \gamma_1 x \sum_{n_1=1}^{\infty} \sum_{n_2=1}^{\infty} n_1 \pi(n_1,n_2) x^{n_1-1} y^{n_2} \\
            & \quad = \quad \gamma_1 x \left[
                \sum_{n_1=1}^{\infty} \sum_{n_2=0}^{\infty} n_1 \pi(n_1,n_2) x^{n_1-1} y^{n_2}
                - \sum_{n_1=1}^{\infty} n_1 \pi(n_1,0) x^{n_1-1}
            \right] \\
            & \quad = \quad \gamma_1 x \left[
                \sum_{n_1=0}^{\infty} \sum_{n_2=0}^{\infty} n_1 \pi(n_1,n_2) x^{n_1-1} y^{n_2}
                - \sum_{n_1=0}^{\infty} n_1 \pi(n_1,0) x^{n_1-1}
            \right] \\
            & \quad = \quad \gamma_1 x \left[
                \frac{\partial}{\partial x} P(x,y) - \frac{d}{dx} P_x(x)
            \right]
        \end{split} \\
        \begin{split}
            & LHS_3 \\
            & \quad = \quad \gamma_2 \sum_{n_1=1}^{\infty} \sum_{n_2=1}^{\infty} n_2 \pi(n_1,n_2) x^{n_1} y^{n_2} \\
            & \quad = \quad \gamma_2 \sum_{n_1=1}^{\infty} \left[ \sum_{n_2=0}^{\infty} n_2 \pi(n_1,n_2) x^{n_1} y^{n_2} - 0 \cdot \pi(n_1,0) x^{n_1} \right] \\
            & \quad = \quad \gamma_2 y \left[
                \sum_{n_1=0}^{\infty} \sum_{n_2=0}^{\infty} n_2 \pi(n_1,n_2) x^{n_1} y^{n_2-1}
                - \sum_{n_2=0}^{\infty} n_2 \pi(0,n_2) y^{n_2-1}
            \right] \\
            & \quad = \quad \gamma_2 y \left[
                \frac{\partial}{\partial y} P(x,y) - \frac{d}{dy} P_y(y)
            \right]
        \end{split} \\
        \begin{split}
            & LHS_4 \\
            & \quad = \quad \theta_1 \sum_{n_1=1}^{\infty} \sum_{n_2=1}^{\infty} n_1 \pi(n_1,n_2) x^{n_1} y^{n_2} \\
            & \quad = \quad \theta_1 x \sum_{n_1=1}^{\infty} \sum_{n_2=1}^{\infty} n_1 \pi(n_1,n_2) x^{n_1-1} y^{n_2} \\
            & \quad = \quad \theta_1 x \left[
                \sum_{n_1=1}^{\infty} \sum_{n_2=0}^{\infty} n_1 \pi(n_1,n_2) x^{n_1-1} y^{n_2}
                - \sum_{n_1=1}^{\infty} n_1 \pi(n_1,0) x^{n_1-1}
            \right] \\
            & \quad = \quad \theta_1 x \left[
                \sum_{n_1=0}^{\infty} \sum_{n_2=0}^{\infty} n_1 \pi(n_1,n_2) x^{n_1-1} y^{n_2}
                - \sum_{n_1=0}^{\infty} n_1 \pi(n_1,0) x^{n_1-1}
            \right] \\
            & \quad = \quad \theta_1 x \left[
                \frac{\partial}{\partial x} P(x,y) - \frac{d}{dx} P_x(x)
            \right]
        \end{split} \\
        \begin{split}
            & LHS_5 \\
            & \quad = \quad \theta_2 \sum_{n_1=1}^{\infty} \sum_{n_2=1}^{\infty} n_2 \pi(n_1,n_2) x^{n_1} y^{n_2} \\
            & \quad = \quad \theta_2 y \left[
                \sum_{n_1=0}^{\infty} \sum_{n_2=0}^{\infty} n_2 \pi(n_1,n_2) x^{n_1} y^{n_2-1}
                - \sum_{n_2=0}^{\infty} n_2 \pi(0,n_2) y^{n_2-1}
            \right] \\
            & \quad = \quad \theta_2 y \left[
                \frac{\partial}{\partial y} P(x,y) - \frac{d}{dy} P_y(y)
            \right]
        \end{split} \\
        \begin{split}
            & RHS_1 \\
            & \quad = \quad \mu \sum_{n_1=1}^{\infty} \sum_{n_2=1}^{\infty} \pi(n_1+1, n_2) x^{n_1} y^{n_2} \\
            & \quad = \quad \mu x^{-1} \sum_{n_1=1}^{\infty} \sum_{n_2=1}^{\infty} \pi(n_1+1, n_2) x^{n_1+1} y^{n_2} \\
            & \quad = \quad \mu x^{-1} \sum_{m=2}^{\infty} \sum_{n_2=1}^{\infty} \pi(m, n_2) x^{m} y^{n_2} \\
            & \quad = \quad \mu x^{-1} \sum_{m=2}^{\infty} \left[
                \sum_{n_2=0}^{\infty} \pi(m, n_2) x^{m} y^{n_2} - \pi(m, 0) x^m
            \right] \\
            & \quad = \quad \mu x^{-1} \left[ \sum_{m=2}^{\infty} \sum_{n_2=0}^{\infty} \pi(m, n_2) x^{m} y^{n_2}
                - \sum_{m=2}^{\infty} \pi(m, 0) x^{m}
            \right] \\
            & \quad = \quad \mu x^{-1} \left[
                \sum_{m=0}^{\infty} \sum_{n_2=0}^{\infty} \pi(m, n_2) x^{m} y^{n_2}
                - x \sum_{n_2=0}^{\infty} \pi(1, n_2) y^{n_2}
                - \sum_{n_2=0}^{\infty} \pi(0, n_2) y^{n_2}
            \right] \\
            & \quad \quad - \; \mu x^{-1} \left[
                \sum_{m=0}^{\infty} \pi(m, 0) x^{m} - x \pi(1, 0) - \pi(0, 0)
            \right] \\
            & \quad = \quad \mu x^{-1} \left[
                P(x,y) - x P_y(y | n_1=1) - P_y(y) - P_x(x) + x \pi(1, 0) + \pi(0, 0)
            \right]
        \end{split} \\
        \begin{split}
            & RHS_2 \\
            & \quad = \quad \lambda_1 \sum_{n_1=1}^{\infty} \sum_{n_2=1}^{\infty} \pi(n_1-1, n_2) x^{n_1} y^{n_2} \\
            & \quad = \quad \lambda_1 x \sum_{n_1=1}^{\infty} \sum_{n_2=1}^{\infty} \pi(n_1-1, n_2) x^{n_1-1} y^{n_2} \\
            & \quad = \quad \lambda_1 x \sum_{m=0}^{\infty} \sum_{n_2=1}^{\infty} \pi(m, n_2) x^{m} y^{n_2} \\
            & \quad = \quad \lambda_1 x \sum_{m=0}^{\infty} \left[
                \sum_{n_2=0}^{\infty} \pi(m, n_2) x^{m} y^{n_2} - \pi(m, 0) x^m
            \right] \\
            & \quad = \quad \lambda_1 x \left[
                \sum_{m=0}^{\infty} \sum_{n_2=0}^{\infty} \pi(m, n_2) x^{m} y^{n_2}
                - \sum_{m=0}^{\infty} \pi(m, 0) x^m
            \right] \\
            & \quad = \quad \lambda_1 x \left[ P(x,y) - P_x(x) \right]
        \end{split} \\
        \begin{split}
            & RHS_3 \\
            & \quad = \quad \lambda_2 \sum_{n_1=1}^{\infty} \sum_{n_2=1}^{\infty} \pi(n_1, n_2-1) x^{n_1} y^{n_2} \\
            & \quad = \quad \lambda_2 y \sum_{n_1=1}^{\infty} \sum_{n_2=1}^{\infty} \pi(n_1, n_2-1) x^{n_1} y^{n_2-1} \\
            & \quad = \quad \lambda_2 y \sum_{n_1=1}^{\infty} \sum_{m=0}^{\infty} \pi(n_1, m) x^{n_1} y^{m} \\
            & \quad = \quad \lambda_2 y \left[
                \sum_{n_1=0}^{\infty} \sum_{m=0}^{\infty} \pi(n_1, m) x^{n_1} y^{m} - \sum_{m=0}^{\infty} \pi(0, m) y^m
            \right] \\
            & \quad = \quad \lambda_2 y \left[ P(x,y) - P_y(y) \right]
        \end{split} \\
        \begin{split}
            & RHS_4 \\
            & \quad = \quad \gamma_1 \sum_{n_1=1}^{\infty} \sum_{n_2=1}^{\infty} (n_1+1) \pi(n_1+1, n_2-1) x^{n_1} y^{n_2} \\
            & \quad = \quad \gamma_1 y \sum_{n_1=1}^{\infty} \sum_{m_2=0}^{\infty} (n_1+1) \pi(n_1+1, m_2) x^{n_1} y^{m_2} \\
            & \quad = \quad \gamma_1 y \sum_{m_1=2}^{\infty} \sum_{m_2=0}^{\infty} m_1 \pi(m_1, m_2) x^{m_1-1} y^{m_2} \\
            & \quad = \quad \gamma_1 y \left[
                \sum_{m_1=0}^{\infty} \sum_{m_2=0}^{\infty} m_1 \pi(m_1, m_2) x^{m_1-1} y^{m_2}
                - \sum_{m_2=0}^{\infty} \pi(1, m_2) y^{m_2}
            \right] \\
            & \quad = \quad \gamma_1 y \left[
                \frac{\partial}{\partial x} P(x,y) - P_y(y | n_1=1)
            \right]
        \end{split} \\
        \begin{split}
            & RHS_5 \\
            & \quad = \quad \gamma_2 \sum_{n_1=1}^{\infty} \sum_{n_2=1}^{\infty} (n_2+1) \pi(n_1-1, n_2+1) x^{n_1} y^{n_2} \\
            & \quad = \quad \gamma_2 \sum_{n_1=1}^{\infty} \sum_{m_2=2}^{\infty} m_2 \pi(n_1-1, m_2) x^{n_1} y^{m_2-1} \\
            & \quad = \quad \gamma_2 x \left[
                \sum_{n_1=1}^{\infty} \sum_{m_2=0}^{\infty} m_2 \pi(n_1-1, m_2) x^{n_1-1} y^{m_2-1}
                - \sum_{n_1=1}^{\infty} \pi(n_1-1, 1) x^{n_1-1}
            \right] \\
            & \quad = \quad \gamma_2 x \left[
                \sum_{m_1=0}^{\infty} \sum_{m_2=0}^{\infty} m_2 \pi(m_1, m_2) x^{m_1} y^{m_2-1}
                - \sum_{m_1=0}^{\infty} \pi(m_1, 1) x^{m_1}
            \right] \\
            & \quad = \quad \gamma_2 x \left[
                \frac{\partial}{\partial y} P(x,y) - P_x(x | n_2=1)
            \right]
        \end{split} \\
        \begin{split}
            & RHS_6 \\
            & \quad = \quad \theta_1 \sum_{n_1=1}^{\infty} \sum_{n_2=1}^{\infty} (n_1+1) \pi(n_1+1, n_2) x^{n_1} y^{n_2} \\
            & \quad = \quad \theta_1 \sum_{m_1=2}^{\infty} \sum_{n_2=1}^{\infty} m_1 \pi(m_1, n_2) x^{m_1-1} y^{n_2} \\
            & \quad = \quad \theta_1 \left[
                \sum_{m_1=0}^{\infty} \sum_{n_2=1}^{\infty} m_1 \pi(m_1, n_2) x^{m_1-1} y^{n_2}
                - \sum_{n_2=1}^{\infty} \pi(1, n_2) y^{n_2}
            \right] \\
            & \quad = \quad \theta_1 \left[
                \sum_{m_1=0}^{\infty} \sum_{n_2=0}^{\infty} m_1 \pi(m_1, n_2) x^{m_1-1} y^{n_2}
                - \sum_{m_1=0}^{\infty} m_1 \pi(m_1, 0) x^{m_1-1}
            \right] \\
            & \quad \quad - \theta_1 \left[
                \sum_{n_2=0}^{\infty} \pi(1, n_2) y^{n_2}
                - \pi(1, 0)
            \right] \\
            & \quad = \quad \theta_1 \left[
                \frac{\partial}{\partial x} P(x,y) - \frac{d}{dx} P_x(x) - P_y(y | n_1=1) + \pi(1, 0)
            \right]
        \end{split} \\
        \begin{split}
            & RHS_7 \\
            & \quad = \quad \theta_2 \sum_{n_1=1}^{\infty} \sum_{n_2=1}^{\infty} (n_2+1) \pi(n_1, n_2+1) x^{n_1} y^{n_2} \\
            & \quad = \quad \theta_2 \sum_{n_1=1}^{\infty} \sum_{m_2=2}^{\infty} m_2 \pi(n_1, m_2) x^{n_1} y^{m_2-1} \\
            & \quad = \quad \theta_2 \left[
                \sum_{n_1=1}^{\infty} \sum_{m_2=0}^{\infty} m_2 \pi(n_1, m_2) x^{n_1} y^{m_2-1}
                - \sum_{n_1=1}^{\infty} \pi(n_1, 1) x^{n_1}
            \right] \\
            & \quad = \quad \theta_2 \left[
                \sum_{n_1=0}^{\infty} \sum_{m_2=0}^{\infty} m_2 \pi(n_1, m_2) x^{n_1} y^{m_2-1}
                - \sum_{m_2=0}^{\infty} m_2 \pi(0, m_2) y^{m_2-1}
            \right] \\
            & \quad \quad - \theta_2 \left[
                \sum_{n_1=0}^{\infty} \pi(n_1, 1) x^{n_1}
                - \pi(0, 1)
            \right] \\
            & \quad = \quad \theta_2 \left[
                \frac{\partial}{\partial y} P(x,y) - \frac{d}{dy} P_y(y) - P_x(x | n_2=1) + \pi(0, 1)
            \right]
        \end{split}
    \end{align}
    
    Now, we combine all terms, we group them and we multiply both sides by $xy$ in order to avoid negative exponents on $x$ and $y$, be it on that or subsequent equations:
    \begin{align}
        \begin{split}
            & (\lambda_1 + \lambda_2 + \mu) \left[
                P(x,y) - P_y(y) - P_x(x) + \pi(0,0)
            \right] \\
            & + \gamma_1 x \left[ \frac{\partial}{\partial x} P(x,y) - \frac{d}{dx} P_x(x) \right] \\
            & + \gamma_2 y \left[ \frac{\partial}{\partial y} P(x,y) - \frac{d}{dy} P_y(y) \right] \\
            & + \theta_1 x \left[ \frac{\partial}{\partial x} P(x,y) - \frac{d}{dx} P_x(x) \right] \\
            & + \theta_2 y \left[ \frac{\partial}{\partial y} P(x,y) - \frac{d}{dy} P_y(y) \right] \\
            & \quad = \quad \mu x^{-1} \left[
                P(x,y) - x P_y(y | n_1=1) - P_y(y) - P_x(x) + x \pi(1,0) + \pi(0,0)
            \right] \\
            & \quad \quad + \lambda_1 x \left[ P(x,y) - P_x(x) \right] \\
            & \quad \quad + \lambda_2 y \left[ P(x,y) - P_y(y) \right] \\
            & \quad \quad + \gamma_1 y \left[ \frac{\partial}{\partial x} P(x,y) - P_y(y | n_1=1) \right] \\
            & \quad \quad + \gamma_2 x \left[ \frac{\partial}{\partial y} P(x,y) - P_x(x | n_2=1) \right] \\
            & \quad \quad + \theta_1 \left[
                \frac{\partial}{\partial x} P(x,y) - \frac{d}{dx} P_x(x) - P_y(y | n_1=1) + \pi(1, 0)
            \right] \\
            & \quad \quad + \theta_2 \left[
                \frac{\partial}{\partial y} P(x,y) - \frac{d}{dy} P_y(y) - P_x(x | n_2=1) + \pi(0, 1)
            \right]
        \end{split} \nonumber \\
        \begin{split}
            & \left[ (\lambda_1+\lambda_2+\mu) - \mu x^{-1} - \lambda_1 x - \lambda_2 y \right] P(x,y) \\
            & + \left[ \gamma_1(x-y) + \theta_1(x-1) \right] \frac{\partial}{\partial x} P(x,y) \\
            & + \left[ \gamma_2(y-x) + \theta_2(y-1) \right] \frac{\partial}{\partial y} P(x,y) \\
            & \quad = \quad \left[ (\lambda_1 + \lambda_2 + \mu) - \mu x^{-1} - \lambda_1 x \right] P_x(x) \\
            & \quad \quad + \left[ (\lambda_1 + \lambda_2 + \mu) - \mu x^{-1} - \lambda_2 y \right] P_y(y) \\
            & \quad \quad + \left[ -(\lambda_1 + \lambda_2 + \mu) + \mu x^{-1} \right] \pi(0,0) \\
            & \quad \quad + \left[ \gamma_1 x + \theta_1(x-1) \right] \frac{d}{dx} P_x(x) \\
            & \quad \quad + \left[ \gamma_2 y + \theta_2(y-1) \right] \frac{d}{dy} P_y(y) \\
            & \quad \quad - \left[ \mu + \gamma_1 y + \theta_1 \right] P_y(y | n_1=1) \\
            & \quad \quad - \left[ \gamma_2 x + \theta_2 \right] P_x(x | n_2=1) \\
            & \quad \quad + \mu \pi(1,0) \\
            & \quad \quad + \theta_1 \pi(1,0) \\
            & \quad \quad + \theta_2 \pi(0,1)
        \end{split} \nonumber \\
        \begin{split}
            & \left[ (\lambda_1+\lambda_2+\mu)xy - \mu y - \lambda_1 x^2 y - \lambda_2 x y^2 \right] P(x,y) \\
            & + xy \left[ \gamma_1(x-y) + \theta_1(x-1) \right] \frac{\partial}{\partial x} P(x,y) \\
            & + xy \left[ \gamma_2(y-x) + \theta_2(y-1) \right] \frac{\partial}{\partial y} P(x,y) \\
            & \quad = \quad \left[ (\lambda_1 + \lambda_2 + \mu)xy - \mu y - \lambda_1 x^2 y \right] P_x(x) \\
            & \quad \quad + \left[ (\lambda_1 + \lambda_2 + \mu)xy - \mu y - \lambda_2 x y^2 \right] P_y(y) \\
            & \quad \quad - \left[ (\lambda_1 + \lambda_2 + \mu)xy - \mu y \right] \pi(0,0) \\
            & \quad \quad + xy \left[ \gamma_1 x + \theta_1(x-1) \right] \frac{d}{dx} P_x(x) \\
            & \quad \quad + xy \left[ \gamma_2 y + \theta_2(y-1) \right] \frac{d}{dy} P_y(y) \\
            & \quad \quad - xy \left[ \mu + \gamma_1 y + \theta_1 \right] P_y(y | n_1=1) \\
            & \quad \quad - xy \left[ \gamma_2 x + \theta_2 \right] P_x(x | n_2=1) \\
            & \quad \quad + (\mu + \theta_1) xy \, \pi(1,0) \\
            & \quad \quad + \theta_2 xy \, \pi(0,1)
        \end{split} \label{eq:gen_transform_interior_relation}
    \end{align}
    
    Now, we can recreate the equivalent steps for the $x$-boundary balance equation (Equation~\eqref{eq:gen_balance_x-boundary}):
    \begin{align}
        \begin{split}
            & \sum_{n_1=1}^{\infty} \left[ \lambda_1 + \lambda_2 + \mu + (\gamma_1 + \theta_1) n_1 \right] \pi(n_1, 0) x^{n_1} \\
            & \quad = \quad \lambda_1 \sum_{n_1=1}^{\infty} \pi(n_1-1,0) x^{n_1}
            + \mu \sum_{n_1=1}^{\infty} \pi(n_1+1,0) x^{n_1} \\
            & \quad \quad + \gamma_2 \sum_{n_1=1}^{\infty} \pi(n_1-1, 1) x^{n_1}
            + \theta_1 \sum_{n_1=1}^{\infty} (n_1+1) \pi(n_1+1, 0) x^{n_1}
            + \theta_2 \sum_{n_1=1}^{\infty} \pi(n_1, 1) x^{n_1}
        \end{split} \nonumber \\
        \begin{split}
            & (\lambda_1 + \lambda_2 + \mu) \sum_{n_1=1}^{\infty} \pi(n_1, 0) x^{n_1}
            + (\gamma_1 + \theta_1) x \sum_{n_1=1}^{\infty} n_1 \pi(n_1, 0) x^{n_1-1} \\
            & \quad = \quad \lambda_1 x \sum_{m=0}^{\infty} \pi(m, 0) x^{m}
            + \mu x^{-1} \sum_{m=2}^{\infty} \pi(m, 0) x^{m} \\
            & \quad \quad + \gamma_2 x \sum_{m=0}^{\infty} \pi(m, 1) x^{m}
            + \theta_1 \sum_{m=2}^{\infty} m \pi(m, 0) x^{m-1}
            + \theta_2 \sum_{m=1}^{\infty} \pi(m, 1) x^{m}
        \end{split} \nonumber \\
        \begin{split}
            & (\lambda_1 + \lambda_2 + \mu) \left[ \sum_{n_1=0}^{\infty} \pi(n_1, 0) x^{n_1} - \pi(0,0) \right]
            + (\gamma_1 + \theta_1) x \sum_{n_1=0}^{\infty} n_1 \pi(n_1, 0) x^{n_1-1} \\
            & \quad = \quad \lambda_1 x \sum_{m=0}^{\infty} \pi(m, 0) x^{m} \\
            & \quad \quad + \mu x^{-1} \left[ \sum_{m=0}^{\infty} \pi(m, 0) x^{m} - x \pi(1, 0) - \pi(0,0) \right] \\
            & \quad \quad + \gamma_2 x \sum_{m=0}^{\infty} \pi(m, 1) x^{m} \\
            & \quad \quad + \theta_1 \left[ \sum_{m=0}^{\infty} m \pi(m, 0) x^{m-1} - \pi(1, 0) \right] \\
            & \quad \quad + \theta_2 \left[ \sum_{m=0}^{\infty} \pi(m, 1) x^{m} - \pi(0, 1) \right]
        \end{split} \nonumber \\
        \begin{split}
            & (\lambda_1 + \lambda_2 + \mu) \left[ P_x(x) - \pi(0,0) \right]
            + (\gamma_1 + \theta_1) x \frac{d}{dx} P_x(x) \\
            & \quad = \quad \lambda_1 x P_x(x)
            + \mu x^{-1} \left[ P_x(x) - x \pi(1, 0) - \pi(0, 0) \right] \\
            & \quad \quad + \gamma_2 x P_x(x | n_2=1)
            + \theta_1 \left[ \frac{d}{dx} P_x(x) - \pi(1, 0) \right]
            + \theta_2 \left[ P_x(x | n_2=1) - \pi(0, 1) \right]
        \end{split} \nonumber \\
        \begin{split}
            & \left[ (\lambda_1 + \lambda_2 + \mu) - \mu x^{-1} - \lambda_1 x \right] P_x(x)
            + \left[ \gamma_1 x + \theta_1(x-1) \right] \frac{d}{dx} P_x(x) \\
            & \quad = \quad \left[ (\lambda_1 + \lambda_2 + \mu) - \mu x^{-1} \right] \pi(0,0)
            - (\mu + \theta_1) \pi(1, 0)
            - \theta_2 \pi(0, 1) \\
            & \quad \quad + \left[ \gamma_2 x + \theta_2 \right] P_x(x | n_2=1)
        \end{split} \nonumber \\
        \begin{split}
            & \left[ (\lambda_1 + \lambda_2 + \mu)xy - \mu y - \lambda_1 x^2 y \right] P_x(x)
            + xy \left[ \gamma_1 x + \theta_1(x-1) \right] \frac{d}{dx} P_x(x) \\
            & \quad = \quad \left[ (\lambda_1 + \lambda_2 + \mu)xy - \mu y \right] \pi(0,0) \\
            & \quad \quad - (\mu + \theta_1) xy \, \pi(1, 0) \\
            & \quad \quad - \theta_2 xy \, \pi(0, 1) \\
            & \quad \quad + xy \left[ \gamma_2 x + \theta_2 \right] P_x(x | n_2=1)
        \end{split} \label{eq:gen_transform_x-boundary_relation}
    \end{align}
    
    Again, we replicate the procedure for the $y$-boundary balance equation (Equation~\eqref{eq:gen_balance_y-boundary}):
    \begin{align}
        \begin{split}
            & \sum_{n_2=1}^{\infty} \left[ \lambda_1 + \lambda_2 + \mu + (\gamma_2 + \theta_2) n_2 \right] \pi(0, n_2) y^{n_2} \\
            & \quad = \quad \lambda_2 \sum_{n_2=1}^{\infty} \pi(0, n_2-1) y^{n_2} \\
            & \quad \quad + \mu \sum_{n_2=1}^{\infty} \pi(1, n_2) y^{n_2} \\
            & \quad \quad + \mu \sum_{n_2=1}^{\infty} \pi(0, n_2+1) y^{n_2} \\
            & \quad \quad + \gamma_1 \sum_{n_2=1}^{\infty} \pi(1, n_2-1) y^{n_2} \\
            & \quad \quad + \theta_1 \sum_{n_2=1}^{\infty} \pi(1, n_2) y^{n_2} \\
            & \quad \quad + \theta_2 \sum_{n_2=1}^{\infty} (n_2+1) \pi(0, n_2+1) y^{n_2}
        \end{split} \nonumber \\
        \begin{split}
            & (\lambda_1 + \lambda_2 + \mu) \sum_{n_2=1}^{\infty} \pi(0, n_2) y^{n_2} \\
            & + (\gamma_2 + \theta_2) y \sum_{n_2=1}^{\infty} n_2 \pi(0, n_2) y^{n_2-1} \\
            & \quad = \quad \lambda_2 y \sum_{m=0}^{\infty} \pi(0, m) y^{m} \\
            & \quad \quad + \mu \sum_{n_2=1}^{\infty} \pi(1, n_2) y^{n_2} \\
            & \quad \quad + \mu y^{-1} \sum_{m=2}^{\infty} \pi(0, m) y^{m} \\
            & \quad \quad + \gamma_1 y \sum_{m=0}^{\infty} \pi(1, m) y^{m} \\
            & \quad \quad + \theta_1 \sum_{n_2=1}^{\infty} \pi(1, n_2) y^{n_2} \\
            & \quad \quad + \theta_2 \sum_{m=2}^{\infty} m \pi(0, m) y^{m-1}
        \end{split} \nonumber \\
        \begin{split}
            & (\lambda_1 + \lambda_2 + \mu) \left[ \sum_{n_2=0}^{\infty} \pi(0, n_2) y^{n_2} - \pi(0,0) \right] \\
            & + (\gamma_2 + \theta_2) y \sum_{n_2=0}^{\infty} n_2 \pi(0, n_2) y^{n_2-1} \\
            & \quad = \quad \lambda_2 y \sum_{m=0}^{\infty} \pi(0, m) y^{m} \\
            & \quad \quad + \mu \left[ \sum_{n_2=0}^{\infty} \pi(1, n_2) y^{n_2} - \pi(1, 0) \right] \\
            & \quad \quad + \mu y^{-1} \left[ \sum_{m=0}^{\infty} \pi(0, m) y^{m} - y \pi(0, 1) - \pi(0, 0) \right] \\
            & \quad \quad + \gamma_1 y \sum_{m=0}^{\infty} \pi(1, m) y^{m} \\
            & \quad \quad + \theta_1 \left[ \sum_{n_2=0}^{\infty} \pi(1, n_2) y^{n_2} - \pi(1, 0) \right] \\
            & \quad \quad + \theta_2 \left[ \sum_{m=0}^{\infty} m \pi(0, m) y^{m-1} - \pi(0, 1) \right]
        \end{split} \nonumber \\
        \begin{split}
            & (\lambda_1 + \lambda_2 + \mu) \left[ P_y(y) - \pi(0, 0) \right]
            + (\gamma_2 + \theta_2) y \frac{d}{dy} P_y(y) \\
            & \quad = \quad \lambda_2 y P_y(y)
            + \mu \left[ P_y(y | n_1=1) - \pi(1, 0) \right] \\
            & \quad \quad + \mu y^{-1} \left[ P_y(y) - y \pi(0, 1) - \pi(0, 0) \right]
            + \gamma_1 y P_y(y | n_1=1) \\
            & \quad \quad + \theta_1 \left[ P_y(y | n_1=1) - \pi(1, 0) \right]
            + \theta_2 \left[ \frac{d}{dy} P_y(y) - \pi(0, 1) \right]
        \end{split} \nonumber \\
        \begin{split}
            & \left[ (\lambda_1 + \lambda_2 + \mu) - \mu y^{-1} - \lambda_2 y \right] P_y(y)
            + \left[ \gamma_2 y + \theta_2(y-1) \right] \frac{d}{dy} P_y(y) \\
            & \quad = \quad \left[ \mu + \gamma_1 y + \theta_1 \right] P_y(y | n_1=1) \\
            & \quad \quad - (\mu + \theta_1) \pi(1, 0)
            - (\mu + \theta_2) \pi(0, 1) \\
            & \quad \quad + \left[ (\lambda_1 + \lambda_2 + \mu) - \mu y^{-1} \right] \pi(0, 0)
        \end{split} \nonumber \\
        \begin{split}
            & \left[ (\lambda_1 + \lambda_2 + \mu) - \mu y^{-1} - \lambda_2 y \right] P_y(y)
            + \left[ \gamma_2 y + \theta_2(y-1) \right] \frac{d}{dy} P_y(y) \\
            & \quad = \quad \left[ \mu + \gamma_1 y + \theta_1 \right] P_y(y | n_1=1) \\
            & \quad \quad + \left[ (\lambda_1 + \lambda_2 + \mu) - \mu y^{-1} - (\lambda_1 + \lambda_2) \right] \pi(0, 0)
        \end{split} \nonumber \\
        \begin{split}
            & \left[ (\lambda_1 + \lambda_2 + \mu) - \mu y^{-1} - \lambda_2 y \right] P_y(y)
            + \left[ \gamma_2 y + \theta_2(y-1) \right] \frac{d}{dy} P_y(y) \\
            & \quad = \quad \left[ \mu + \gamma_1 y + \theta_1 \right] P_y(y | n_1=1)
            + \left[ \mu - \mu y^{-1} \right] \pi(0, 0)
        \end{split} \nonumber \\
        \begin{split}
            & \left[ (\lambda_1 + \lambda_2 + \mu)xy - \mu x - \lambda_2 x y^2 \right] P_y(y)
            + xy \left[ \gamma_2 y + \theta_2(y-1) \right] \frac{d}{dy} P_y(y) \\
            & \quad = \quad \left[ \mu + \gamma_1 y + \theta_1 \right] xy \, P_y(y | n_1=1)
            - \mu x(1-y) \, \pi(0,0)
        \end{split} \label{eq:gen_transform_y-boundary_relation}
    \end{align}
    
    We can now combine Equation~\eqref{eq:gen_transform_interior_relation} and Equation~\eqref{eq:gen_transform_x-boundary_relation}, where many terms cancel out, leaving us with the following, remarkably cleaner, expression:
    \begin{align}
        \begin{split}
            & \left[ (\lambda_1+\lambda_2+\mu)xy - \mu y - \lambda_1 x^2 y - \lambda_2 x y^2 \right] P(x,y) \\
            & + xy \left[ \gamma_1(x-y) + \theta_1(x-1) \right] \frac{\partial}{\partial x} P(x,y) \\
            & + xy \left[ \gamma_2(y-x) + \theta_2(y-1) \right] \frac{\partial}{\partial y} P(x,y) \\
            & \quad = \quad \left[ (\lambda_1 + \lambda_2 + \mu)xy - \mu y - \lambda_2 x y^2 \right] P_y(y) \\
            & \quad \quad + xy \left[ \gamma_2 y + \theta_2(y-1) \right] \frac{d}{dy} P_y(y) \\
            & \quad \quad - xy \left[ \mu + \gamma_1 y + \theta_1 \right] P_y(y | n_1=1)
        \end{split} \label{eq:gen_plugged_interior_and_x}
    \end{align}
    
    Now, isolating $\left[ \mu + \gamma_1 y + \theta_1 \right] xy \, P_y(y | n_1=1)$ in Equation~\eqref{eq:gen_transform_y-boundary_relation} and plugging it into Equation~\eqref{eq:gen_plugged_interior_and_x} makes a cascade of cancelations that result in:
    \begin{align}
        \begin{split}
            & \left[ (\lambda_1+\lambda_2+\mu)xy - \mu y - \lambda_1 x^2 y - \lambda_2 x y^2 \right] P(x,y) \\
            & + xy \left[ \gamma_1(x-y) + \theta_1(x-1) \right] \frac{\partial}{\partial x} P(x,y) \\
            & + xy \left[ \gamma_2(y-x) + \theta_2(y-1) \right] \frac{\partial}{\partial y} P(x,y) \\
            & \quad = \quad \left[ (\lambda_1 + \lambda_2 + \mu)xy - \mu y - \lambda_2 x y^2 \right] P_y(y) \\
            & \quad \quad + xy \left[ \gamma_2 y + \theta_2(y-1) \right] \frac{d}{dy} P_y(y) \\
            & \quad \quad - \left\{ \left[ (\lambda_1 + \lambda_2 + \mu)xy - \mu x - \lambda_2 x y^2 \right] P_y(y) \right. \\
            & \quad \quad \quad \left. + xy \left[ \gamma_2 y + \theta_2(y-1) \right] \frac{d}{dy} P_y(y) + \mu x(1-y) \, \pi(0,0) \right\}
        \end{split} \nonumber \\
        \begin{split}
            & \left[ (\lambda_1+\lambda_2+\mu)xy - \mu y - \lambda_1 x^2 y - \lambda_2 x y^2 \right] P(x,y) \\
            & + xy \left[ \gamma_1(x-y) + \theta_1(x-1) \right] \frac{\partial}{\partial x} P(x,y) \\
            & + xy \left[ \gamma_2(y-x) + \theta_2(y-1) \right] \frac{\partial}{\partial y} P(x,y) \\
            & \quad = \quad \mu(x-y) P_y(y) - \mu x(1-y) \, \pi(0,0)
        \end{split} \label{eq:gen_fundamental_equation}.
    \end{align}
\end{proof}

Sanity check $1$, $\gamma_1=\gamma_2=\theta_1=\theta_2=0$ (Model A):
\begin{equation}
\begin{split}
    & \left[ (\lambda_1+\lambda_2+\mu)xy - \mu y - \lambda_1 x^2 y - \lambda_2 x y^2 \right] P(x,y) \\
    & \quad = \quad \mu(x-y) P_y(y) - \mu x(1-y) \, \pi(0,0)
\end{split}
\label{eq:n1n2_sanity_check_1}
\end{equation}

Sanity check $2$, $\gamma_1=\gamma_2=0$ (abandonments only):
\begin{equation}
\begin{split}
    & \left[ (\lambda_1+\lambda_2+\mu)xy - \mu y - \lambda_1 x^2 y - \lambda_2 x y^2 \right] P(x,y) \\
    & + \theta_1 xy (x-1) \frac{\partial}{\partial x} P(x,y) \\
    & + \theta_2 xy (y-1) \frac{\partial}{\partial y} P(x,y) \\
    & \quad = \quad \mu(x-y) P_y(y) - \mu x(1-y) \, \pi(0,0)
\end{split}
\label{eq:n1n2_sanity_check_2}
\end{equation}

Sanity check $3$, $\theta_1=\theta_2=0$, recovers Model B:
\begin{equation}
\begin{split}
    & \left[ (\lambda_1+\lambda_2+\mu)xy - \mu y - \lambda_1 x^2 y - \lambda_2 x y^2 \right] P(x,y) \\
    & + \gamma_1 xy (x-y) \frac{\partial}{\partial x} P(x,y) \\
    & + \gamma_2 xy (y-x) \frac{\partial}{\partial y} P(x,y) \\
    & \quad = \quad \mu(x-y) P_y(y) - \mu x(1-y) \, \pi(0,0)
\end{split}
\label{eq:n1n2_sanity_check_3}
\end{equation}


\subsubsection{Equation for PPGF in state space $\widetilde{S}$}

\begin{lemma}[Fundamental equation for the PPGF on $\widetilde{S}$]\label{lem:gen_PPGF_dynamics}
    Consider the same system on the state space
    $\widetilde{S}=\{(n_2,n):0\le n_2\le n\}$, where $n$ is the total number of customers
    \emph{waiting} in the queues and $n_2$ the number of waiting class-$2$ customers.
    Assume positive recurrence, so that the stationary distribution
    $\{\tpi(n_2,n)\}$ exists, and define the partial PGF (transforming only $n_2$,
    keeping $n$ discrete)
    \[
        \widetilde{P}(y,n)=\sum_{n_2=0}^{n}\tpi(n_2,n)\,y^{n_2},
    \]
    a polynomial of degree $n$ in $y$. Then, for every $n\ge 1$, the family
    $\{\widetilde{P}(y,n)\}_{n\ge 0}$ satisfies the differential--difference equation
    \begin{equation}
    \begin{split}
        & \left[ -(\gamma_2 + \gamma_1 y)(1-y) + (\theta_2 - \theta_1) y \right] \frac{d}{dy} \widetilde{P}(y,n)
          + \left[ \lambda_1 + \lambda_2 + \mu + \gamma_1 n(1-y) + \theta_1 n \right] \widetilde{P}(y,n) \\
        & \quad = \quad (\lambda_1 + \lambda_2 y) \widetilde{P}(y,n-1)
          + \left[ \mu + \theta_1 (n+1) \right] \widetilde{P}(y,n+1) \\
        & \quad \quad + (\theta_2 - \theta_1 y) \frac{d}{dy} \widetilde{P}(y,n+1)
          + \mu y^{n} (1-y)\, \tpi(n+1,n+1).
    \end{split}
    \label{eq:gen_PPGF_dynamics}
    \end{equation}
\end{lemma}

\begin{proof}
    We start by taking the interior balance equation (Equation~\eqref{eq:tilde_balance_interior}), multiplying by $y^{n_2}$ and adding up throughout its domain (i.e., $1 \leq n_2 \leq n-1$):
    \begin{align*}
        \begin{split}
            & (\lambda_1+\lambda_2+\mu) \sum_{n_2=1}^{n-1} \tpi(n_2,n) y^{n_2} \\
            & + \gamma_1 \sum_{n_2=1}^{n-1} (n-n_2) \tpi(n_2,n) y^{n_2} \\
            & + \gamma_2 \sum_{n_2=1}^{n-1} n_2 \tpi(n_2,n) y^{n_2} \\
            & + \theta_1 \sum_{n_2=1}^{n-1} (n-n_2) \tpi(n_2,n) y^{n_2} \\
            & + \theta_2 \sum_{n_2=1}^{n-1} n_2 \tpi(n_2,n) y^{n_2} \\
            & \quad = \quad \lambda_1 \sum_{n_2=1}^{n-1} \tpi(n_2, n-1) y^{n_2} \\
            & \quad \quad + \lambda_2 \sum_{n_2=1}^{n-1} \tpi(n_2-1, n-1) y^{n_2} \\
            & \quad \quad + \mu \sum_{n_2=1}^{n-1} \tpi(n_2, n+1) y^{n_2} \\
            & \quad \quad + \gamma_1 \sum_{n_2=1}^{n-1} (n-n_2+1) \tpi(n_2-1, n) y^{n_2} \\
            & \quad \quad + \gamma_2 \sum_{n_2=1}^{n-1} (n_2+1) \tpi(n_2+1, n) y^{n_2} \\
            & \quad \quad + \theta_1 \sum_{n_2=1}^{n-1} (n+1-n_2) \tpi(n_2, n+1) y^{n_2} \\
            & \quad \quad + \theta_2 \sum_{n_2=1}^{n-1} (n_2+1) \tpi(n_2+1, n+1) y^{n_2}
        \end{split}
    \end{align*}
    
    We add the diagonal balance equation (Equation~\eqref{eq:tilde_balance_diagonal}) multiplied by $y^n$ for the $n$-th term.
    
    When adding these two equations, we notice that most terms from the diagonal equation are cleanly absorbed by the interior balance equation, extending the limit of the summation up to $n$. On the LHS, this is the case for the $(\lambda_1+\lambda_2+\mu)$-term, the $\gamma_2$-term and the $\theta_2$-term; the $\gamma_1$-term and the $\theta_1$-term are extended automatically since their $n$-th term carries the factor $(n-n_2)$ which vanishes at $n_2 = n$. On the RHS, the $\lambda_2$-term, the $\mu$-term, the $\gamma_1$-term, the $\theta_1$-term and the $\theta_2$-term are absorbed by the diagonal contributions. This leaves all summations completed excepting the $\lambda_1$-term and the $\gamma_2$-term on the RHS, plus a spare $\mu \tpi(n+1,n+1) y^n$ also on the RHS:
    \begin{align*}
        \begin{split}
            & (\lambda_1+\lambda_2+\mu) \sum_{n_2=1}^{n} \tpi(n_2,n) y^{n_2} \\
            & + \gamma_1 \sum_{n_2=1}^{n} (n-n_2) \tpi(n_2,n) y^{n_2} \\
            & + \gamma_2 \sum_{n_2=1}^{n} n_2 \tpi(n_2,n) y^{n_2} \\
            & + \theta_1 \sum_{n_2=1}^{n} (n-n_2) \tpi(n_2,n) y^{n_2} \\
            & + \theta_2 \sum_{n_2=1}^{n} n_2 \tpi(n_2,n) y^{n_2} \\
            & \quad = \quad \lambda_1 \sum_{n_2=1}^{n-1} \tpi(n_2,n-1) y^{n_2} \\
            & \quad \quad + \lambda_2 \sum_{n_2=1}^{n} \tpi(n_2-1,n-1) y^{n_2} \\
            & \quad \quad + \mu \sum_{n_2=1}^{n} \tpi(n_2,n+1) y^{n_2} \\
            & \quad \quad + \mu \tpi(n+1,n+1) y^n \\
            & \quad \quad + \gamma_1 \sum_{n_2=1}^{n} (n-n_2+1) \tpi(n_2-1,n) y^{n_2} \\
            & \quad \quad + \gamma_2 \sum_{n_2=1}^{n-1} (n_2+1) \tpi(n_2+1,n) y^{n_2} \\
            & \quad \quad + \theta_1 \sum_{n_2=1}^{n} (n+1-n_2) \tpi(n_2,n+1) y^{n_2} \\
            & \quad \quad + \theta_2 \sum_{n_2=1}^{n} (n_2+1) \tpi(n_2+1,n+1) y^{n_2}.
        \end{split}
    \end{align*}
    
    Now we work the terms out individually:
    \begin{align}
        \begin{split}
            & LHS_1 \\
            & \quad = \quad (\lambda_1 + \lambda_2 + \mu) \sum_{n_2=1}^{n} \tpi(n_2,n) y^{n_2} \\
            & \quad = \quad (\lambda_1 + \lambda_2 + \mu) \left[ \sum_{n_2=0}^{n} \tpi(n_2,n) y^{n_2} - \tpi(0,n) y^0 \right] \\
            & \quad = \quad (\lambda_1 + \lambda_2 + \mu) \left[ \widetilde{P}(y,n) - \tpi(0,n) \right]
        \end{split} \\
        \begin{split}
            & LHS_2 \\
            & \quad = \quad \gamma_1 \sum_{n_2=1}^{n} (n-n_2) \tpi(n_2,n) y^{n_2} \\
            & \quad = \quad \gamma_1 \left[ n \sum_{n_2=1}^{n} \tpi(n_2,n) y^{n_2} - \sum_{n_2=1}^{n} n_2 \tpi(n_2,n) y^{n_2} \right] \\
            & \quad = \quad \gamma_1 \left[ n \left( \sum_{n_2=0}^{n} \tpi(n_2,n) y^{n_2} - \tpi(0,n) y^0 \right) - y \sum_{n_2=1}^{n} n_2 \tpi(n_2,n) y^{n_2-1} \right] \\
            & \quad = \quad \gamma_1 \left[ n \big( \widetilde{P}(y,n) - \tpi(0,n) \big) - y \frac{d}{dy} \widetilde{P}(y,n) \right]
        \end{split} \\
        \begin{split}
            & LHS_3 \\
            & \quad = \quad \gamma_2 \sum_{n_2=1}^{n} n_2 \tpi(n_2,n) y^{n_2} \\
            & \quad = \quad \gamma_2 \left[ y \sum_{n_2=1}^{n} n_2 \tpi(n_2,n) y^{n_2-1} \right] \\
            & \quad = \quad \gamma_2 y \frac{d}{dy} \widetilde{P}(y,n)
        \end{split} \\
        \begin{split}
            & LHS_4 \\
            & \quad = \quad \theta_1 \sum_{n_2=1}^{n} (n-n_2) \tpi(n_2,n) y^{n_2} \\
            & \quad = \quad \theta_1 \left[ n \big( \widetilde{P}(y,n) - \tpi(0,n) \big) - y \frac{d}{dy} \widetilde{P}(y,n) \right]
        \end{split} \\
        \begin{split}
            & LHS_5 \\
            & \quad = \quad \theta_2 \sum_{n_2=1}^{n} n_2 \tpi(n_2,n) y^{n_2} \\
            & \quad = \quad \theta_2 y \frac{d}{dy} \widetilde{P}(y,n)
        \end{split} \\
        \begin{split}
            & RHS_1 \\
            & \quad = \quad \lambda_1 \sum_{n_2=1}^{n-1} \tpi(n_2,n-1) y^{n_2} \\
            & \quad = \quad \lambda_1 \left[ \sum_{n_2=1}^{n} \tpi(n_2,n-1) y^{n_2} - \tpi(n,n-1) y^n \right] \\
            & \quad = \quad \lambda_1 \left[ \sum_{n_2=0}^{n} \tpi(n_2,n-1) y^{n_2} - \tpi(n,n-1) y^n - \tpi(0,n-1) y^0 \right] \\
            & \quad = \quad \lambda_1 \left[ \widetilde{P}(y,n-1) - \tpi(n,n-1) y^n - \tpi(0,n-1) \right]
        \end{split} \\
        \begin{split}
            & RHS_2 \\
            & \quad = \quad \lambda_2 \sum_{n_2=1}^{n} \tpi(n_2-1,n-1) y^{n_2} \\
            & \quad = \quad \lambda_2 \left[ y \sum_{n_2=1}^{n} \tpi(n_2-1,n-1) y^{n_2-1} \right] \\
            & \quad = \quad \lambda_2 \left[ y \sum_{n_2=0}^{n-1} \tpi(n_2,n-1) y^{n_2} \right] \\
            & \quad = \quad \lambda_2 \left[ y \left( \sum_{n_2=0}^{n} \tpi(n_2,n-1) y^{n_2} - \tpi(n,n-1) y^n \right) \right] \\
            & \quad = \quad \lambda_2 y \left[ \widetilde{P}(y,n-1) - \tpi(n,n-1) y^n \right]
        \end{split} \\
        \begin{split}
            & RHS_3 \\
            & \quad = \quad \mu \tpi(n+1,n+1) y^n + \mu \sum_{n_2=1}^{n} \tpi(n_2,n+1) y^{n_2} \\
            & \quad = \quad \mu \left[ \tpi(n+1,n+1) y^n + \sum_{n_2=1}^{n} \tpi(n_2,n+1) y^{n_2} \right] \\
            & \quad = \quad \mu \left[ \tpi(n+1,n+1) y^n + \sum_{n_2=0}^{n+1} \tpi(n_2,n+1) y^{n_2} - \tpi(0,n+1) y^0 - \tpi(n+1,n+1) y^{n+1} \right] \\
            & \quad = \quad \mu \left[ \widetilde{P}(y,n+1) - \tpi(0,n+1) + \tpi(n+1,n+1) y^n (1-y) \right]
        \end{split} \\
        \begin{split}
            & RHS_4 \\
            & \quad = \quad \gamma_1 \sum_{n_2=1}^{n} (n-n_2+1) \tpi(n_2-1,n) y^{n_2} \\
            & \quad = \quad \gamma_1 \left[ y \sum_{n_2=1}^{n} (n - (n_2-1)) \tpi(n_2-1,n) y^{n_2-1} \right] \\
            & \quad = \quad \gamma_1 \left[ y \sum_{k=0}^{n-1} (n-k) \tpi(k,n) y^k \right] \\
            & \quad = \quad \gamma_1 \left[ y \sum_{k=0}^{n} (n-k) \tpi(k,n) y^k \right] \quad \text{(since the $k=n$ term is 0)} \\
            & \quad = \quad \gamma_1 y \left[ n \sum_{k=0}^{n} \tpi(k,n) y^k - y \sum_{k=1}^{n} k \tpi(k,n) y^{k-1} \right] \\
            & \quad = \quad \gamma_1 y \left[ n \widetilde{P}(y,n) - y \frac{d}{dy} \widetilde{P}(y,n) \right]
        \end{split} \\
        \begin{split}
            & RHS_5 \\
            & \quad = \quad \gamma_2 \sum_{n_2=1}^{n-1} (n_2+1) \tpi(n_2+1,n) y^{n_2} \\
            & \quad = \quad \gamma_2 \sum_{k=2}^{n} k \tpi(k,n) y^{k-1} \quad \text{(letting $k = n_2+1$)} \\
            & \quad = \quad \gamma_2 \left[ \sum_{k=1}^{n} k \tpi(k,n) y^{k-1} - 1 \cdot \tpi(1,n) y^0 \right] \\
            & \quad = \quad \gamma_2 \left[ \frac{d}{dy} \widetilde{P}(y,n) - \tpi(1,n) \right]
        \end{split} \\
        \begin{split}
            & RHS_6 \\
            & \quad = \quad \theta_1 \sum_{n_2=1}^{n} (n+1-n_2) \tpi(n_2,n+1) y^{n_2} \\
            & \quad = \quad \theta_1 \left[ (n+1) \sum_{n_2=1}^{n} \tpi(n_2,n+1) y^{n_2} - \sum_{n_2=1}^{n} n_2 \tpi(n_2,n+1) y^{n_2} \right] \\
            & \quad = \quad \theta_1 \left[ (n+1) \left( \sum_{n_2=0}^{n+1} \tpi(n_2,n+1) y^{n_2} - \tpi(0,n+1) - \tpi(n+1,n+1) y^{n+1} \right) \right] \\
            & \quad \quad - \theta_1 \left[ y \left( \sum_{n_2=0}^{n+1} n_2 \tpi(n_2,n+1) y^{n_2-1} - (n+1) \tpi(n+1,n+1) y^n \right) \right] \\
            & \quad = \quad \theta_1 \left[ (n+1) \big( \widetilde{P}(y,n+1) - \tpi(0,n+1) \big) - (n+1) \tpi(n+1,n+1) y^{n+1} \right] \\
            & \quad \quad - \theta_1 \left[ y \frac{d}{dy} \widetilde{P}(y,n+1) - (n+1) \tpi(n+1,n+1) y^{n+1} \right] \\
            & \quad = \quad \theta_1 \left[ (n+1) \big( \widetilde{P}(y,n+1) - \tpi(0,n+1) \big) - y \frac{d}{dy} \widetilde{P}(y,n+1) \right]
        \end{split} \\
        \begin{split}
            & RHS_7 \\
            & \quad = \quad \theta_2 \sum_{n_2=1}^{n} (n_2+1) \tpi(n_2+1,n+1) y^{n_2} \\
            & \quad = \quad \theta_2 \sum_{k=2}^{n+1} k \tpi(k,n+1) y^{k-1} \quad \text{(letting $k = n_2+1$)} \\
            & \quad = \quad \theta_2 \left[ \sum_{k=1}^{n+1} k \tpi(k,n+1) y^{k-1} - 1 \cdot \tpi(1,n+1) y^0 \right] \\
            & \quad = \quad \theta_2 \left[ \frac{d}{dy} \widetilde{P}(y,n+1) - \tpi(1,n+1) \right]
        \end{split}
    \end{align}
    
    Put the components back together:
    \begin{align*}
        \begin{split}
            & (\lambda_1 + \lambda_2 + \mu) \left[ \widetilde{P}(y,n) - \tpi(0,n) \right] \\
            & + \gamma_1 \left[ n \big( \widetilde{P}(y,n) - \tpi(0,n) \big) - y \frac{d}{dy} \widetilde{P}(y,n) \right] \\
            & + \gamma_2 y \frac{d}{dy} \widetilde{P}(y,n) \\
            & + \theta_1 \left[ n \big( \widetilde{P}(y,n) - \tpi(0,n) \big) - y \frac{d}{dy} \widetilde{P}(y,n) \right] \\
            & + \theta_2 y \frac{d}{dy} \widetilde{P}(y,n) \\
            & \quad = \quad \lambda_1 \left[ \widetilde{P}(y,n-1) - \tpi(n,n-1) y^n - \tpi(0,n-1) \right] \\
            & \quad \quad + \lambda_2 y \left[ \widetilde{P}(y,n-1) - \tpi(n,n-1) y^n \right] \\
            & \quad \quad + \mu \left[ \widetilde{P}(y,n+1) - \tpi(0,n+1) + \tpi(n+1,n+1) y^n (1-y) \right] \\
            & \quad \quad + \gamma_1 y \left[ n \widetilde{P}(y,n) - y \frac{d}{dy} \widetilde{P}(y,n) \right] \\
            & \quad \quad + \gamma_2 \left[ \frac{d}{dy} \widetilde{P}(y,n) - \tpi(1,n) \right] \\
            & \quad \quad + \theta_1 \left[ (n+1) \big( \widetilde{P}(y,n+1) - \tpi(0,n+1) \big) - y \frac{d}{dy} \widetilde{P}(y,n+1) \right] \\
            & \quad \quad + \theta_2 \left[ \frac{d}{dy} \widetilde{P}(y,n+1) - \tpi(1,n+1) \right]
        \end{split}
    \end{align*}
    Grouping terms:
    \begin{align*}
        \begin{split}
            & \left[ -(\gamma_2 + \gamma_1 y)(1-y) + (\theta_2 - \theta_1) y \right] \frac{d}{dy} \widetilde{P}(y,n) + \left[ \lambda_1 + \lambda_2 + \mu + \gamma_1 n(1-y) + \theta_1 n \right] \widetilde{P}(y,n) \\
            & \quad = \quad (\lambda_1 + \lambda_2 y) \widetilde{P}(y,n-1) + \left[ \mu + \theta_1 (n+1) \right] \widetilde{P}(y,n+1) + (\theta_2 - \theta_1 y) \frac{d}{dy} \widetilde{P}(y,n+1) \\
            & \quad \quad + \left\{ \left[ \lambda_1 + \lambda_2 + \mu + (\gamma_1 + \theta_1) n \right] \tpi(0,n) - \lambda_1 \tpi(0,n-1) - \left[ \mu + \theta_1 (n+1) \right] \tpi(0,n+1) \right. \\
            & \quad \quad \quad \left. - \gamma_2 \tpi(1,n) - \theta_2 \tpi(1,n+1) \right\} \\
            & \quad \quad - y^n (\lambda_1 + \lambda_2 y) \tpi(n,n-1) + \mu y^n (1-y) \tpi(n+1,n+1)
        \end{split}
    \end{align*}
    Combine with the $y$-boundary balance equation (Equation~\eqref{eq:tilde_balance_y-boundary}), which makes the bracketed boundary terms vanish. Together with $\tpi(n,n-1)=0$ (state outside $\widetilde{S}$), the equation simplifies to:
    \begin{align}
        \begin{split}
            & \left[ -(\gamma_2 + \gamma_1 y)(1-y) + (\theta_2 - \theta_1) y \right] \frac{d}{dy} \widetilde{P}(y,n) + \left[ \lambda_1 + \lambda_2 + \mu + \gamma_1 n(1-y) + \theta_1 n \right] \widetilde{P}(y,n) \\
            & \quad = \quad (\lambda_1 + \lambda_2 y) \widetilde{P}(y,n-1) + \left[ \mu + \theta_1 (n+1) \right] \widetilde{P}(y,n+1) + (\theta_2 - \theta_1 y) \frac{d}{dy} \widetilde{P}(y,n+1) \\
            & \quad \quad + \mu y^n (1-y) \tpi(n+1,n+1)
        \end{split}
    \end{align}
\end{proof}

Sanity check $1$, $\gamma_1=\gamma_2=\theta_1=\theta_2=0$ (Model A):
\begin{align}
    \begin{split}
        & (\lambda_1 + \lambda_2 + \mu) \widetilde{P}(y,n) \\
        & \quad = \quad (\lambda_1 + \lambda_2 y) \widetilde{P}(y,n-1) + \mu \widetilde{P}(y,n+1) + \mu y^n (1-y) \tpi(n+1,n+1)
    \end{split}
\end{align}

Sanity check $2$, $\theta_1 = \theta_2 = 0$ (recovers Model B):
\begin{align}
    \begin{split}
        & -(\gamma_2 + \gamma_1 y)(1-y) \frac{d}{dy} \widetilde{P}(y,n) + \left[ \lambda_1 + \lambda_2 + \mu + \gamma_1 n(1-y) \right] \widetilde{P}(y,n) \\
        & \quad = \quad (\lambda_1 + \lambda_2 y) \widetilde{P}(y,n-1) + \mu \widetilde{P}(y,n+1) + \mu y^n (1-y) \tpi(n+1,n+1)
    \end{split}
\end{align}

Sanity check $3$, $\gamma_1 = \gamma_2 = 0$ (abandonments only):
\begin{align}
    \begin{split}
        & (\theta_2 - \theta_1) y \frac{d}{dy} \widetilde{P}(y,n) + \left[ \lambda_1 + \lambda_2 + \mu + \theta_1 n \right] \widetilde{P}(y,n) \\
        & \quad = \quad (\lambda_1 + \lambda_2 y) \widetilde{P}(y,n-1) + \left[ \mu + \theta_1 (n+1) \right] \widetilde{P}(y,n+1) \\
        & \quad \quad + (\theta_2 - \theta_1 y) \frac{d}{dy} \widetilde{P}(y,n+1) + \mu y^n (1-y) \tpi(n+1,n+1)
    \end{split}
\end{align}


